% TOP_PROBLEM: {"id": 30006797, "title": "Wall's Poincare Duality Group Conjecture", "category_id": 17, "category": {"id": 17, "name": "group_theory", "display_name": "Group Theory", "description": "Problems about groups, group actions, representations, and related algebraic structures.", "slug": "group-theory", "order_index": 17, "created_at": "2026-07-31T15:26:25.671Z"}, "status": "open"}
% FIRST_POSTED: 2026-09-25
% ENTRY_KIND: statement_only
% !TeX program = lualatex
% Definition and sources only; this file is not an LLM solution attempt.
\documentclass[11pt]{article}
\usepackage[margin=25mm]{geometry}
\usepackage{fontspec}
\setmainfont{Latin Modern Roman}
\usepackage{amsmath,amssymb,mathtools}
\usepackage{unicode-math}
\setmathfont{Latin Modern Math}
\usepackage{xurl,hyperref}
\hypersetup{colorlinks=true,urlcolor=blue}
\setcounter{secnumdepth}{0}
\setlength{\parindent}{0pt}
\setlength{\parskip}{5pt}
\setlength{\emergencystretch}{3em}
\begin{document}
\section*{\texorpdfstring{Wall's Poincare Duality Group Conjecture}{Wall's Poincare Duality Group Conjecture}}\label{30006797}
\subsection{Definitions and mathematical statement}
An integral Poincar\'e duality group of dimension $n$ is a group of cohomological dimension $n$, of type FP over ${\mathbb{Z}}$, whose cohomology with group-ring coefficients vanishes outside degree $n$ and is infinite cyclic in degree $n$. The action on that cyclic dualizing module defines its orientation character $w:G\to\{\pm1\}$. For every finitely presented such $G$ and every $n\ge3$, ask for a closed aspherical topological $n$-manifold $M$ with
\[
 \pi_1(M)\cong G,
\]
inducing the same orientation character. Aspherical means the universal cover is contractible. A finite Poincar\'e duality complex is not automatically a topological manifold, and that distinction is central to the target.
\par\smallskip\textit{Formulation and concept references:} \href{https://arxiv.org/pdf/0902.2480}{[S1]}.

\subsection{Short English statement}
Is every finitely presented integral Poincaré duality group of dimension at least three the fundamental group of a closed aspherical manifold of that dimension?
\subsection{Sources}
\begin{itemize}
\item[S1]Wolfgang Lück, “Survey on Aspherical Manifolds,” arXiv:0902.2480v3 (2009); published in European Mathematical Society proceedings (2010).
\url{https://arxiv.org/pdf/0902.2480}.
\textit{Formulation source.}
\end{itemize}
\end{document}
